\subsection{Macaulay 代数或 Gorenstein 代数的张量积}

命题 4. 设 k 为域，A 为本质有限生成 k-代数，B 为 k-代数。

\begin{enumerate}
\def\labelenumi{\alph{enumi})}
\item
  若 A 和 B 是 Macaulay 环，则 A ⊗k B 亦然。
\item
  若 A 和 B 是 Gorenstein 环，则 \(\Lambda\otimes_{k}\mathrm{B}\) 亦然。
\end{enumerate}

设 A 和 B 是 Macaulay（相应地 Gorenstein）环，我们来证明 \(A \otimes_{k} B\) 亦然。环 \(A \otimes_{k} B\) 是 Noether 的 (第 1 节, 命题 2 的推论)。A-模 \(A \otimes_{k} B\) 是自由的，从而是平坦的。由 § 2, 第 7 节的命题 10（相应地 § 3, 第 8 节的命题 12 的推论 1），只需证明：对 A 的每个素理想 p，\(\kappa(\mathfrak{p}) \otimes_{k} B\) 都是 Macaulay 环（相应地 Gorenstein 环）。由于扩张 \(\kappa(\mathfrak{p})\) 在 k 上是有限型的 (第 1 节, 命题 2 及例 1)，我们于是归结为证明 k-代数 A 是 k 的有限生成扩张 K 这一情形的命题。

设 \((t_{1},\ldots,t_{n})\) 为 K 在 k 上的一个超越基，于是 K 是纯扩张 \(k^{\prime}=k(t_{1},\ldots,t_{n})\) 的有限次扩张 (A, V, 第 112 页, 命题 17)。环 \(B^{\prime}=k^{\prime}\otimes_{k}B\) 同构于多项式环 \(B[T_{1},\ldots,T_{n}]\) 的一个分式环，从而由 § 2, 第 7 节的命题 10 的推论 2（相应地 § 3, 第 8 节的命题 12 的推论 3）是 Macaulay（相应地 Gorenstein）环。由于 \(K\otimes_{k}B\) 等同于 \(K\otimes_{k^{\prime}}B^{\prime}\)，我们归结为证明 K 是 k 的有限次扩张时断言 a) 和 b)。

环 K ⊗k B 此时是一个有限秩自由 B-模，从而是 Macaulay B-模。由 § 2, 第 6 节的命题 8 的推论 1，它是 Macaulay 环，这就证明了 a)。现在设 B 是 Gorenstein 环。存在 K 的子扩张 K \(_{i}\)（0 ≤ i ≤ m）：

\[
k = \mathrm{K} _ {0} \subset \mathrm{K} _ {1} \subset \dots \subset \mathrm{K} _ {m} = \mathrm{K}
\]

使 \(K_{i}\) 是 \(K_{i-1}\) -单生成代数（对 \(i = 1, \ldots, m\)），而这使我们能归结到 k 的扩张 K 是（有限次的）单生成的情形。典范同态 \(B \to K \otimes_{k} B\) 使 \(K \otimes_{k} B\) 成为平坦 B-代数，且对每个 \(\mathfrak{q} \in \operatorname{Spec}(B)\)，环 \((\mathrm{K} \otimes_{k} \mathrm{B}) \otimes_{\mathrm{B}} \mathfrak{k}(\mathfrak{q})\)（同构于 \(K \otimes_{k} \mathfrak{k}(\mathfrak{q})\)）是 \(\mathfrak{k}(\mathfrak{q})\) -单生成代数，从而是 Gorenstein 环 (§ 3, 第 9 节, 例 5)。于是 \(K \otimes_{k} B\) 是 Gorenstein 环 (§ 3, 第 8 节, 命题 12 的推论 1)，这就完成了 b) 的证明。

推论 1.--- 设 \(k\) 为域，K 为 \(k\) 的扩张，\(\Lambda\) 为本质有限型 \(k\) -代数。为使 \(\mathrm{A}_{(\mathrm{K})}\) 是 Macaulay 环（相应地 Gorenstein 环），必须且只需 A 亦然。

若 A 是 Macaulay（相应地 Gorenstein）环，则由命题 4，\(\mathrm{A}_{(\mathrm{K})}\) 亦然。由于 \(\mathrm{A}_{(\mathrm{K})}\) 是忠实平坦 A-模，逆命题由 § 2, 第 7 节的命题 10（相应地 § 3, 第 8 节的命题 12 的推论 1）得出。

推论 2.--- 设 k 为 Noether 环，A 和 B 为 k-代数。假设 A 在 k 上平坦且本质有限型。若 A 和 B 是 Macaulay 环（相应地 Gorenstein 环），则 A ⊗k B 是 Macaulay 环（相应地 Gorenstein 环）。

先设 B 是域；用 \(\varphi\) 表示从 \(k\) 到 B 中的典范同态，用 \(\mathfrak{r}\) 表示它的核。同态 \(\varphi\) 诱导一个从分式域 \(\kappa(\mathfrak{r})\)（\(k/\mathfrak{r}\) 的）到 B 中的同态。于是 A \(\otimes_{k}\) B 等同于 (A \(\otimes_{k}\) K(r)) \(\otimes_{\kappa(\mathfrak{r})}\) B；由于 \(\varphi^{-1}(0)=\mathfrak{r}\)，环 A \(\otimes_{k}\) K(r) 由 § 2, 第 7 节的命题 10（相应地 § 3, 第 8 节的命题 12 的推论 1）是 Macaulay（相应地 Gorenstein）环。此时断言由推论 1 得出。

现在转到一般情形。B-代数 A⊗ \(_{k}\) B 是平坦的，且由第 1 节的命题 2 的推论是 Noether 的。对 B 的每个素理想 q，环 (A⊗ \(_{k}\) B)⊗ \(_{B}\) κ(q)（它等同于 A⊗ \(_{k}\) κ(q)）由前面的结果是 Macaulay（相应地 Gorenstein）环。由 § 2, 第 7 节的命题 10（相应地 § 3, 第 8 节的命题 12 的推论 1）即可得出结论。

